\chapter{Glossary}
\label{chap:glossary}

\begin{description}[leftmargin=*,style=nextline]
    \item[Accident year]
        The calendar year in which a claim occurred.

    \item[Aggregate loss]
        The total loss across all claims in a simulation period.

    \item[Aggregate Exceedance Probability (AEP)]
        The probability that aggregate (annual) losses exceed a given
        threshold. Reported as a function of return period in
        catastrophe risk applications.

    \item[Bornhuetter--Ferguson method]
        A reserving method that combines an a priori expected ultimate
        loss ratio with reported development. Available in the
        \texttt{chainladder} package.

    \item[Chain-ladder method]
        A standard reserving technique that projects ultimate losses
        from a triangle by applying age-to-age development factors.

    \item[Claim development]
        The evolution of a claim's reported or paid amount between its
        date of occurrence and ultimate settlement.

    \item[Copula]
        A multivariate distribution function with uniform marginals,
        used to introduce dependency between random variables while
        preserving their marginal distributions.

    \item[Cumulative development factor (CDF)]
        The ratio of ultimate losses to losses at a given development
        age. The product of LDFs from a given age through to ultimate.

    \item[Development age]
        The number of months elapsed since the start of an accident
        period.

    \item[Frequency]
        The number of claims occurring during a specified period.

    \item[Goodness-of-fit metric]
        A statistic used to compare candidate distributions, such as
        AIC, BIC, log-likelihood, chi-square, or Kolmogorov--Smirnov.

    \item[Loss development factor (LDF)]
        The ratio of losses at one development age to losses at the
        previous age.

    \item[Monte Carlo simulation]
        A simulation technique that estimates the distribution of an
        outcome by repeated random sampling.

    \item[Non-homogeneous Poisson process (NHPP)]
        A Poisson process with a time-varying intensity function.
        Used in \actsim to assign occurrence dates to simulated
        claims with seasonal patterns.

    \item[Occurrence Exceedance Probability (OEP)]
        The probability that the largest single-event loss in a year
        exceeds a given threshold.

    \item[Paid loss]
        The cumulative amount paid on a claim at a given valuation
        date.

    \item[Random seed]
        An integer that initialises a pseudo-random number generator.
        Setting the seed makes simulations reproducible.

    \item[Reported loss]
        The sum of paid losses and case reserves on a claim at a given
        valuation date.

    \item[Severity]
        The amount of loss associated with an individual claim.

    \item[Tail Value at Risk (TVaR)]
        The expected loss conditional on exceeding the VaR threshold,
        also known as the conditional tail expectation.

    \item[Triangle]
        A standard actuarial data layout in which rows correspond to
        accident periods and columns correspond to development ages.

    \item[Ultimate loss]
        The expected total cost of a claim once all development is
        complete.

    \item[Value at Risk (VaR)]
        The quantile of a loss distribution at a given confidence level.
\end{description}
